\section{Discussion}\label{sec:discussion}

\subsection{Principal anticipated findings}\label{subsec:principal-findings}

The principal anticipated finding of the four-month deployment is empirical evidence that the four-modality acquisition pipeline (UCLM ad-hoc ambient platform, EmotiBit physiological wristbands, Bitbrain Versatile 32-channel semi-dry EEG and overhead 4K body-pose features) can be operated simultaneously in ordinary 3rd- and 4th-of-ESO Mathematics teaching, under the European data-protection regime for minors, over a sustained longitudinal window. The validation block reports the per-modality measurement properties under realistic deployment conditions; the descriptive characterisation of the four-month corpus reports the marginal and joint distributions of the per-window teacher label, the monthly PVT change scores and the per-modality feature aggregates. The anticipated novelty is not in any single modality --- each is individually mature --- but in the conjunction: an integrated, synchronised, longitudinal, ecologically valid, externally anchored, privacy-preserving deployment in compulsory secondary education that meets the four conjunctive conditions documented in \Cref{sec:introduction}. The release of the protocol \emph{before} any predictive modelling allows the multimodal learning analytics community to scrutinise the measurement decisions before downstream inferential claims are made.

\subsection{Anticipated downstream applications: multimodal AI fusion and teacher feedback}

\textcolor{blue}{While the primary scope of this protocol is observational data acquisition, the resulting curated dataset is designed to support downstream multimodal machine learning models. The synchronised environmental, physiological, EEG and body-pose streams provide a foundation for training classifiers that disambiguate student states --- for example, distinguishing elevated physiological arousal driven by poor air quality from arousal linked to cognitive demand \parencite{di_lascio_unobtrusive_2018,vanneste_towards_2021}. A teacher-facing dashboard, such as the proof-of-concept interface illustrated in Figure~\ref{fig:dashboard}, could render aggregated environmental and physiological trends to support situational awareness during instrumented sessions. These downstream applications are deferred to subsequent compendium articles; the present protocol commits only to the measurement infrastructure that makes them possible.}

\begin{figure}[H]
    \centering
    \includegraphics[width=0.9\textwidth]{figuras/dashboard.png}
    \caption{\textbf{Proof-of-concept teacher dashboard.} The interface renders aggregated environmental parameters (temperature, equivalent CO$_2$, humidity, illuminance, sound pressure level) alongside physiological telemetry from the wearable biosensors (heart rate, electrodermal activity, accelerometry). The dashboard is a development prototype for future descriptive analytics; no real-time feedback is delivered to teachers during the protocol deployment.}
    \label{fig:dashboard}
\end{figure}

\subsection{Strengths}\label{subsec:strengths}

The protocol's strengths follow from the design discipline imposed at the protocol level rather than from any single technical choice. First, the simultaneous integration of four passive modalities under a single NTP-disciplined acquisition pipeline is, to the team's knowledge, the first such deployment in compulsory secondary education at the longitudinal window described above. Second, the 5-minute binary teacher annotation provides an interpretable ground-truth label that is generable in situ by the same teachers who will use the downstream predictive system, while the monthly Psychomotor Vigilance Task supplies an independent external benchmark of sustained attention whose sensitivity and convergent validity are well established \parencite{Dinges1985_pvt,Basner2011_pvt,Loh2004_pvt_short}. Third, the privacy-by-design regime is operationally enforceable: raw classroom video is deleted within 72 hours of body-pose feature extraction under a verifiable governance log, audio is never captured at any device, and raw EEG and EmotiBit traces remain on the institutional UCLM server (\Cref{sec:appendix-privacy}). The arrangement implements the GDPR Article 9 / LOPDGDD requirements for special-category data of minors and the data-protection-by-design principles of Article 25 \parencite{Macenaite2017_minors_gdpr,EDPB2020_consent}. Fourth, the protocol is the measurement foundation of a doctoral compendium: subsequent articles will analyse multimodal predictive performance, the relative contribution of each modality, the minimal sensor subset that retains predictive utility and the moderating role of the teaching context, all on the dataset and feature definitions documented here.

\subsection{Comparison with prior work}\label{subsec:comparison}

The protocol sits in a literature whose state of the art the team surveyed systematically. Individual deployments meet subsets of the conjunctive conditions documented in \Cref{sec:introduction}: classroom-scale computer vision pipelines for engagement and attention have been validated under naturalistic conditions \parencite{Hutt2019_mind_wandering,Parambil2022_smart_classroom}; multimodal sensing platforms for higher-education contexts have demonstrated the integration of vision, audio and physiological streams in instrumented classrooms \parencite{Ahuja2019_edusense}; multimodal learning analytics has matured as a conceptual programme \parencite{Blikstein2016_multimodal,Schneider2015_unraveling}. The present protocol differs from these precedents in the conjunction --- ambient + physiological + EEG + body-pose, longitudinal over four months, compulsory secondary education, PVT-anchored, GDPR Article 9 for minors --- and in the explicit commitment to characterise the platform's own measurement properties before any predictive claim is made.

\subsection{Construct validity and convergent validation}\label{subsec:construct-validity}

The protocol uses cognitive performance as an umbrella construct that is operationalised through three distinct measurements: the teacher's binary judgement of predominantly on-task behaviour per 5-minute window, the participant-level Psychomotor Vigilance Task administered monthly, and an end-of-session 0--10 global rating. On-task behaviour, teacher impression of attention and psychomotor vigilance are not the same construct in the educational-psychology and operational-cognition literatures. They are treated jointly here as convergent indicators of cognitive performance under the explicit assumption that systematic departures between them are diagnostic --- not noise --- and will themselves form part of the protocol's descriptive output. Specifically, the participant-session-level correlation between the within-session $\textrm{óptimo}$-rate and the PVT change score (PVT pre-session minus PVT post-session, oriented so that improvement is positive) is the convergent-validity test that the team commits to report; the absence of correlation would be an empirical finding about the operational definitions, not a failure of the platform.

\subsection{Falsifiability thresholds}\label{subsec:falsifiability}

The protocol commits in advance to two falsification thresholds. First, if the per-session yield falls below 70\% of expected samples on two or more modalities for more than 25\% of recording weeks, the team will issue a corrigendum stating that the protocol as published does not, in fact, deliver usable multimodal data under the described conditions and will redesign before publishing any predictive analysis. Second, if the participant-level Spearman correlation between the PVT change score and the within-session $\textrm{óptimo}$-rate is statistically indistinguishable from zero across the four months (95\% bootstrap confidence interval containing zero), the team will publish a substantive note reconsidering the convergent-validity assumption between the two operationalisations of cognitive performance, including a re-examination of whether the binary scheme should be replaced with a multi-dimensional engagement coding in subsequent compendium articles. These thresholds are documented here so that the protocol's outcome is not retro-interpretable.

\subsection{Limitations}\label{subsec:limitations}

The protocol's limitations follow from the design choices made to keep the deployment operationally viable and ethically defensible in a four-month classroom field study with minors. \textbf{Single school, single subject, single region.} The deployment is one secondary school in Albacete during ordinary Mathematics sessions of one 3rd-of-ESO group, with the comparison group recorded during Mathematics sessions of the same school year. External validity to other schools, subjects or regions is not characterised by the present design and will require replication; the article does not make a claim of geographical or subject-area generalisation. \textbf{Teacher anecdotal record as ground truth.} The 5-minute binary label is captured by a single annotator, the teacher leading the session, through Google Forms with timed prompts; a single-annotator design introduces annotator-level label noise and possible drift across sessions \parencite{Frenay2014_label_noise,Bosch2022_engagement_observers}, and the inter-rater reliability that a multi-annotator design would provide is not assessed. A spot-check Cohen's $\kappa$ within-teacher between prompt-time annotation and post-session reconstruction is the residual descriptor of label reliability committed in \Cref{subsec:analysis-plan}. \textcolor{blue}{\textbf{Teacher annotation fatigue.} The annotation load on a single teacher---one binary judgement every 5 minutes across four sessions per week for four months---is substantial. Annotation quality may drift with fatigue, routine, or changes in the teacher's own attentional state. The within-teacher spot-check analysis described above is designed to detect gross drift; subtler degradation would require a second annotator, which the single-teacher design does not provide.} \textbf{Single EEG participant per week.} The Bitbrain Versatile 32-channel system is deployed on one stratified-sampled participant per week, yielding approximately 16 individual EEG sessions over four months; statistical power for participant-level EEG analyses is bounded by design, and EEG inferences are interpreted at the group level.

\textbf{Observational, no experimental control.} By design, the study is purely observational; the causal claims pursued at the primary-education level by the parent national project PID2022-137397NB-I00 on the effect of active breaks on cognitive performance are not transferable to this thesis, and the analyses that the dataset supports are descriptive and predictive rather than interventional. \textbf{Wear acceptance and signal quality in adolescents.} Adolescent participants may exhibit higher motion contamination and lower wear compliance than the university or adult populations on which the wearable devices have been most extensively validated; the protocol mitigates this through trained personnel for EEG fitting and a non-dominant-forearm wear pattern for EmotiBit, but residual quality loss is expected and will be characterised quantitatively \parencite{LauZhu2019_mobile_eeg_review}. \textbf{Practice effects on the external benchmark.} Even at the validated 5-minute format the PVT produces small practice effects under monthly administration \parencite{Basner2018_pvt_practice}; the protocol accommodates this through change scores relative to the participant's pre-session baseline rather than absolute scores. \textbf{Visibility of participation in a class-complete deployment.} The visibility of EmotiBit wristbands and the per-week EEG headset may, in principle, function as a behavioural spotlight; the mitigations are detailed in \Cref{subsec:eeg-spotlight}, but residual social-comparison risk is acknowledged. \textcolor{blue}{\textbf{Monitoring reactivity.} The sustained presence of sensors in the classroom---environmental mote, overhead camera, wristbands, and the rotating EEG headset---may alter student and teacher behaviour beyond the initial familiarisation period. A four-month deployment attenuates novelty effects, but the protocol cannot fully disentangle true cognitive-performance variation from reactive behavioural changes induced by the monitoring itself. The comparison group, which is not instrumented with sensors, provides a partial control for this effect, although between-group differences in monitoring reactivity cannot be excluded.}

\subsection{Conclusions}\label{subsec:conclusions}

The protocol describes an observational, longitudinal multimodal passive sensing study in Spanish compulsory secondary education classrooms, designed to characterise its own measurement properties before any inferential modelling and to support the subsequent predictive analyses of the doctoral compendium. The conjunctive contribution --- four passive modalities, four months, ordinary teaching, PVT external anchor, GDPR Article 9 for minors, privacy-preserving on-server processing, tiered data release --- addresses the structural gap identified in \Cref{sec:introduction}. The explicit pre-commitment to four falsifiable hypotheses and to a residual-reliability descriptor of the primary label ensures that the protocol's outcome is empirically accountable and that the subsequent predictive claims of the compendium rest on a characterised, validated, and transparently documented measurement foundation.
